% ============================================================
%  H. Høffding, "Filosofien og Darwinismen"
%  Et Foredrag, holdt i Studenterforeningen den 21. Marts 1874
%  Nær og Fjern, nr. 93–94, 1874
%  TRANSLATION (English)
% ============================================================
\documentclass[12pt,a4paper]{article}
\usepackage[utf8]{inputenc}
\usepackage[T1]{fontenc}
\usepackage[english]{babel}
\usepackage{microtype}
\usepackage{setspace}
\usepackage[margin=1.2in]{geometry}
\usepackage{fancyhdr}
\usepackage{hyperref}

\hypersetup{
  pdftitle={Philosophy and Darwinism},
  pdfauthor={H. Høffding},
  hidelinks
}

\pagestyle{fancy}
\fancyhf{}
\rhead{Høffding, \emph{Philosophy and Darwinism} (1874)}
\cfoot{\thepage}

\setstretch{1.3}

\title{\textbf{Philosophy and Darwinism}\\[0.4em]
\large A Lecture Delivered at the Students' Association, 21 March}
\author{Dr.\ phil.\ Harald Høffding\\[4pt]
  \small Translated from the Danish}
\date{\emph{Nær og Fjern}, nos.\ 93--94, 1874}

\begin{document}
\maketitle
\thispagestyle{fancy}

\bigskip

It is a familiar fact of experience that oppositions set reflection in
motion. What one previously apprehended with indifference and without
attention suddenly becomes interesting once one catches sight of the
oppositions bound up with it, and of the tensions to which it gives
rise. Someone who, for other reasons, would feel no inclination to
object to Darwinism may nevertheless find a motive for doing so in
the conflicting interpretations of it---the diametrically opposed
judgments of its worth and significance that one encounters on a
daily basis. For one person, Darwinism is a wild speculation, a
product of our materialistic age, which seeks confirmation for its
strange and ultimate drives and inclinations in self-generated
hypotheses; for another, it is a brilliant and magnificent
conception, opening up vast perspectives, satisfying the human
mind's ideal longing for the infinite, and awakening enthusiasm for
higher and progressive forms of development. And if we return to
the narrower circle in which this doctrine arose, and where
scientific testing properly belongs---to the natural scientists
themselves---we likewise find a striking opposition of views. While
a large number of natural scientists (especially English and German
ones) regard Darwin's doctrine as so highly probable that they are
inclined to overlook what is lacking in the evidential support---indeed,
that several even embark on far-reaching speculations and abandon the
firm ground of their discipline---others (notably French researchers,
with whom most Danes seem to side) insist that an explanation is not
a proof, but presupposes proof in order to have validity, and that,
despite all the importance hypotheses may have for science, it is
entirely inadmissible, even for a moment, to forget the gaps in the
evidence because of the appeal of a hypothetical explanation.

There is therefore good reason to undertake an examination of the
logical and ethical worth of Darwinism, in order to settle which of
these competing judgments is the true one, and what stance we ought
to adopt toward this remarkable doctrine, which has now made its way
from the scholar's quiet study out into the streets and is being
employed as a powerful weapon in the intellectual battles of our age.
Philosophy has, apart from these external occasions, by virtue of its
own task, a duty to carry out such an examination. It is philosophy's
task to attain self-consciousness with respect to the nature and
method of scientific cognition, to sharpen insight into the stages a
doctrine must pass through before it can be incorporated into the
realm of scientific truths, and to uphold the necessary conditions for
the truth and validity of cognition. But its vocation is not confined
to this. As it investigates the life of human consciousness, its laws
and its elements, and traces its development through history, it seeks
at the same time to situate every cognitive result as a member within
the whole organism of mind, and must therefore at every point examine
how the result achieved stands as a contribution to a general
worldview. No result has truly been appropriated by the human mind
until this examination has been carried out. And only when this has
been done can the result have legitimate practical consequences, for
the third member of the philosophical investigation becomes: what
practical and ethical significance the appropriated truth acquires.

It is clear that the logical examination is here the most important.
One must first settle whether a doctrine is true or not, before
proceeding to examine its significance in a metaphysical and ethical
respect. One is all too often inclined to take the reverse path and
suppose that what has ideal significance must also have real existence.
From this inclination have sprung the theological and speculative
doctrines that hold sway over the general consciousness. And many of
Darwinism's champions have adopted the same standpoint; they have
fashioned a new theology, a new dogma, out of the purely empirical
hypothesis Darwin put forward. The theory of cognition must oppose
both of these tendencies and seek above all to fix its gaze on the
real connections of the matter. It impresses upon us the necessity of
resignation. An old proverb says: ``One who believes does not hurry.''
One might with equal justice say: ``One who knows does not hurry!''---
understanding by the one who knows the person who has an eye for the
essential conditions of cognition as well as for the ideas that lead
the human mind forward in its theoretical and practical striving. Why
should these ideas necessarily have a---supernatural or natural---
reality? Are they not precisely ideas by virtue of pointing beyond
what is factually given, by posing new tasks for thinking and action?
One whose worldview does not require a dogmatic fixing of the content
of the highest ideas as external fact will, with greater calm and
impartiality, allow the scientific method for the cognition of the
real to be satisfied, and will bow entirely to its results.

\section*{I.}

One may say that science would stand still without hypothesis. In
order to find the law according to which a given factor operates, or
to find the cause underlying a given regularity in phenomena, one
must necessarily venture beyond what is immediately at hand, and
must often make large leaps into entirely different circles of
phenomena in order to test whether the sought explanation is not to
be found there. It has rightly been said that researchers who stand
still and wait until they are driven from outside to move forward
will never make the great discoveries. Discovery presupposes an
activity that goes beyond the mere accumulation of facts, that seeks
a unity in the manifold, a common law for all the differences. It is
not least this activity and energy that makes Copernicus, Kepler, and
Newton such great figures in the history of science. But this is only
one side of the matter. The path thought takes from phenomena to the
law or cause may be the one that looks boldest and most magnificent;
but it is not the one that is, scientifically speaking, decisive. We
may admire that energy, but we do not rely on what it accomplishes
when it cannot manage to travel the opposite path back again, when it
cannot, from the assumed law and the postulated cause, lead us back
to the phenomenal connection. In other words: the hypothesis must be
verified, confirmed by its agreement with facts won through immediate
observation. Newton too had an eye for this, when he demands that the
cause one assumes to explain certain phenomena must be a \textit{vera
causa}---that is, one that in reality proves to have the effects
attributed to it.

What is it, then, that Darwin's hypothesis seeks to explain? It is
the fact that living beings divide themselves into certain distinct
groups with distinctive characteristic marks that remain fixed and
persistent for each individual group---that is, the fact of what are
called specific differences. The plant and animal kingdoms divide
themselves, for observing and classifying inquiry, into sharply
demarcated species. Botanists and zoologists have, on the common view,
completed their task once they have, by careful examination of a
plant's or animal's inner and outer characteristics, determined which
species, genus, and family it must be assigned to. One traces the
developmental history of the individual, one investigates the
agreements among individual species and groups them accordingly; but
the question of the origin of the species itself is not drawn into the
investigation, because immediate observation shows no transition from
one species to another.

Yet there are many things that can indicate that the concept of
species does not have absolute validity. Not only does it prove
difficult, indeed impossible in many cases, to determine what should
be called a species and what should be regarded as a mere variety;
but this very representation of absolutely completed life-forms---
according to which life was only the self-preservation of these
forms, or their struggle to maintain themselves against external
conditions---the representation by which a sharp distinction was
drawn between the organism and its surroundings, as if they did not
necessarily belong together as members of the circle of vital
phenomena, was too contradictory and outwardly mechanical for many
researchers not to feel the urge to go beyond it. So long as
theological ideas still dominated the consciousness of most natural
scientists, this urge could admittedly be satisfied by tracing the
absolute specific differences back to different acts of creation,
each species being conceived as created on its own. But when
awareness of natural connection and the necessity of natural causes
develops, the urge can no longer be satisfied in this way. Insofar as
natural scientists still regard the different species as each
separately created, this is in general only to be taken as a
rhetorical expression of their assumption that the origin of species
is inexplicable. Already at the close of the last century the idea of
the natural origin of species makes its appearance, and remarkably
enough at roughly the same time in England, France, and Germany: in
Erasmus Darwin (Charles Darwin's grandfather), in Geoffroy St.\
Hilaire and Lamarck, and in Goethe and Oken. These earliest attempts
each point, individually, to the operation of external causes, but
most of them lay the main weight on an inner formative force, a
spontaneous capacity for perfection that was supposed to lead living
beings from lower to higher levels. To that extent there is still
something mystical and indeterminate in these attempts.

The brilliance of the way in which Darwin took up the old problem
rests on the fact that this indeterminate and mystical element steps
aside in favour of reference to a definite, factually operative force.
To that extent Newton's demand for a \textit{vera causa} can be said
to be satisfied. The factor Darwin refers to is the so-called quality
selection, or natural selection---the fact that only the favourably
placed organisms live and reproduce, while the weak and unfavourably
placed perish. Nature thus passes living beings through a sieve, as it
were, and from this process of sifting there must necessarily emerge
more developed forms better suited to the conditions. It is this
natural selection that produces races and varieties: why, then, should
it not also be able to produce species? ``It has often been
asserted,'' says Darwin, ``but the assertion is incapable of proof,
that the amount of variation under nature is a strictly limited
quantity\ldots\ I can see no limit to this power, in slowly and
beautifully adapting each form to the most complex relations of life.''
In these words the logical validity of Darwinism finds its clearest
expression.

The resolution of the dispute depends, namely, on the question: does
natural selection have limits or not? But this question science is
unable to answer at its present stage. This is what makes Darwinism a
hypothesis and justifies the hesitation of sober-minded researchers to
see it presented as an established doctrine that one may now simply
accept. This explains the stance of French positivism toward this
question. There are undoubtedly many who conflate positivism and
Darwinism without knowing that they exclude one another. Auguste Comte
sharply criticized Lamarck's theory, seeing in it an unjustified
attempt to transfer the law of development, which applies to the
individual organism, to the species as a whole. One of his disciples,
the celebrated anatomist Charles Robin, has recently repeated this
objection and shown that it applies to Darwin as well. He maintains in
particular, drawing on the exact results of anatomy, that variation
has its definite limits. Every organism can swing between its extreme
points; if the alteration goes beyond them, the organism does not
develop but dies. And precisely the living beings whose organization
is very simple---such as zoophytes and worms---vary only within narrow
limits; yet it is from them that the higher species are supposed to
have developed. And in any case one cannot trace the development of
one species from another in the way one can trace how one stage of an
individual's life develops out of another. ``It is also progress and
a sign of insight to understand how to examine a general proposition
before accepting it.'' Robin accordingly regards it as a testimony to
the philosophical spirit of French researchers that most of them
declare themselves against Darwin. The German Häckel will admittedly
make it a test of human intellectual superiority whether one accepts
Darwinism and ``the philosophy grounded upon it'' or not, and
accordingly credits the English and Germans with being further advanced
in cultural development. Littré has rightly branded this as
chauvinism. And if---which from the standpoint of rigorous science
must be regarded as possible---the Darwinian hypothesis should prove
false, the standing of that superiority would look precarious. What
Kant doubted---that there would ever arise a Newton for the science of
organic life---Häckel sees realized in Darwin. He holds that Darwin's
theory is one of the greatest conquests of the human mind, which
``can immediately be placed alongside Newton's theory of gravitation''
---indeed, that it even surpasses it! Against such an oversight of the
principles of the inductive method as is here laid at the foundation,
logic must protest. Newton's discovery stands as an ideal in science
because it is an unparalleled example of the union of brilliant
intuition and rigorous demonstration, in which no link has been
skipped and all particulars fit together into a magnificent edifice.
Darwin has demonstrated a real cause that is of essential and perhaps
hitherto overlooked significance for organic life, but he has not been
able to prove that this cause contains the sufficient ground for the
origin of the different species.

Will it now be difficult and long before a decisive answer is found to
the question whether variation has definite limits or not? Darwin
himself has set out on the path that alone can lead to an exact proof:
the path of experiment. As is well known, it was domestic animals and
cultivated plants to which Darwin turned his attention above all, and
from which he drew his most important material. If one could produce
sufficiently varied conditions, and for a sufficiently long time, the
question might perhaps find its answer; but it is not at all certain
that it lies within human power to bring about the conditions under
which in free nature species must have developed from one another.
This is therefore only a very indeterminate and remote possibility.

An essential---and, in Darwin's own judgment, the most essential---
objection to his hypothesis is drawn from the fact that ``although the
geological investigations have undoubtedly taught us the existence of
many previous links, which bring numerous life-forms much closer to
one another, they have not given the infinitely many fine gradations
that the theory requires, between extinct and now-living species.''
Here too an exact proof is a long way off---indeed, it is perhaps
doubtful whether such a proof can ever be attained by this route,
since every new intermediate link that is discovered can always be
regarded as a new species to be inserted between two others; the very
transition from one species to another cannot by the nature of the
matter be found in the geological strata, which contain only the dead
remains.

The probability of the hypothesis's truth will of course increase the
more such intermediate links are demonstrated. One also sees that
researchers who were previously opponents of the mutability of species
give up their opposition under the impression of what new discoveries
bring to light. ``That a natural scientist who has for so long and so
thoroughly studied the numerous remains and compared them so carefully
and tirelessly, both with related remains from other places and with
the skeletons of living species---such as Gaudry---has lost his faith
in the absolute fixity of species and aligns himself with those who
believe in their transformation over time, is under all circumstances
a significant pointer, all the more so since he can hardly have been
influenced in his judgment by the natural scientists in whose
proximity he worked; for nowhere has Darwinism won less entry than in
France.''

But an exact proof, as said, does not appear to be within prospect,
at least for a very long time. To that extent those who, for moral
and religious reasons, believe they must combat Darwinism are right
when they mockingly refer to the moment ``when Darwinism will be
proved.'' Yet their satisfaction may nonetheless be somewhat
premature. For it is one thing that the factors by which one has
hitherto sought to explain the origin of species are insufficient; it
is quite another that the origin of species is at all regarded as a
question belonging to natural science and not to theology. And this
latter seems in any case to have been achieved through the Darwinian
controversy. That species must have a natural origin, that the causes
of their arising, whatever they may otherwise be, must be sought in
the conditions given in nature---this is a presupposition that springs
from the principles of all scientific inquiry, principles that are
increasingly being applied even to domains from which they were
previously excluded. The maxim that everything has its cause springs
from the innermost law of human consciousness itself and cannot be
forced back once it has been acknowledged. What the cause is, is
another question. Many hypotheses of a kind similar to Darwin's may
perhaps succeed one another before the true one is found. ``It is,''
says Mill, ``no valid reason for accepting a given hypothesis that we
are unable to imagine any other that can explain the facts. There is
no necessity for supposing that the true explanation must be one which,
with our present experience, we are able to conceive.'' But in
principle it remains fixed that there must be a natural cause. The
skepticism of French positivism is only justified with respect to
individual explanatory attempts, not with regard to the principle
itself. It may be premature to wish to resolve this question already
now, but it does not belong to those that in themselves lie beyond
the limits of science.

\begin{center}
[\emph{To be concluded.}]
\end{center}

% ============================================================
%  SECOND PART (Nær og Fjern, no. 94)
% ============================================================

\bigskip
\begin{center}
\textbf{Second Part}
\end{center}
\medskip

\noindent \emph{Concluded.}

\section*{II.}

The logical examination of Darwinism thus leads to assigning it its
rightful place as a brilliant hypothesis that has grown out of the
principle of modern thought and science---the principle of the
dominion of natural laws and natural causes in existence. This
principle did not wait for Darwinism in order to come into force; it
has found its logical and metaphysical grounding in modern philosophy
and its real implementation in modern natural science. It therefore
does not stand or fall with Darwinism either. We cannot accordingly
attribute to Darwinism so great a significance for the overall
worldview as is often done; it is a valuable, if hypothetical, member
of a larger connection.

Yet it can be said to have become, as it were, a watchword in the
struggle between the supernaturalist and spiritualist conception and
the humane and rational worldview. The objections directed against
Darwinism from that conception would be directed against any other
doctrine that sought to explain the origin of species in a natural
way. It will therefore, entirely apart from Darwinism's hypothetical
character, be of interest to examine what direction it points in when
one draws its final consequences.

The common charge is that Darwinism is a materialist doctrine that
explains the inner from the outer, the spiritual from the corporeal,
and conceives of the human being as an animal by letting him
``descend from the apes.'' In place of the divine wisdom that,
according to the old faith, guides everything in nature and in history
to the good, one sees blind chance holding sway---the apparently
purposive owing its existence only to a fortunate accident, or rather
a series of fortunate accidents.

We take no account here of Darwin's personal views, even though they
do not appear to tend in a materialist direction. In order to test the
charge just stated, we hold to the cause from which he seeks to explain
the origin of species: the struggle for existence. It is, first of
all, clear that struggle presupposes a struggling being, and weapons
and force in this struggling being. That is, there is presupposed an
inner activity in the organism and an original endowment from which
development can take its beginning. Without this, external influences
accomplish nothing. Furthermore, the expression ``struggle for
existence'' also contains the implication that existence is a good
that the individual strives to preserve. For it is evidently not
merely being as such but the enjoyment connected with existence and
the value thereby accruing to it that is the object of the struggle.
But these presuppositions do not point in a materialist direction.
They lead on the contrary to the acknowledgment of ideal principles
as what is truly decisive. Darwin himself has nonetheless probably
laid too great a weight on external causes. The celebrated botanist
Nägeli has, in opposition to this, stated as the result of his
investigations that the formation of varieties and races is not a
consequence and expression of external influences but is conditioned
by inner causes; he accordingly attributes to the organism a
``principle of perfection'' that steadily leads to new forms---a
process in which the external conditions play only the role of a
regulator. We are thereby led back to the general law of life, as
formulated most sharply among the more recent philosophers by Herbert
Spencer: an accommodation of the inner to the outer. The inner is
thus always presupposed as an original factor that, through its
interaction with the surrounding world, works its way forward to more
perfect forms. According to Spencer, Darwin has laid too little weight
on the direct accommodation that occurs through the influence of
function on structure. Only when the organism is unable to accommodate
itself to its surroundings and thereby overcome them does the indirect
accommodation come into force, through the elimination of the weak.

What, then, is the worldview that would have to be derived from
Darwinism, if it proved to be correct? None other than the one that
follows from the whole philosophical and scientific consciousness of
the modern age, which cannot admit of ideal powers that would operate
in any way other than through the real conditions themselves and the
laws that govern them. We can, as a witty thinker has said, believe
in operative, but not in governing Ideas. The ideal is not in itself
alien to the real conditions; these are elements of its being, and it
demonstrates precisely its divinity by working its way forward through
all obstacles. The Idea we believe in must be a struggling Idea.

It is so far from being the case that Darwinism leads to materialism
that it must rather be said to point in the wholly opposite direction.
It does certainly lay a strong emphasis on external conditions; but
these conditions operate only through the organism's accommodation to
them. And what one calls chance is only an instance of the infinitely
ramified and involved connection of conditions---inner and outer
conditions---that we cannot survey; in what are called contingently
arisen forms we have precisely a testimony to the energy of life,
which knows how to appropriate even the finest nuances of external
conditions and employ them in its service.

But, it is said, Darwin's doctrine is comfortless, because it shows
us that the weak perish and only the strong and fortunate survive. How
dreadful to know oneself and everything else delivered over to such a
dark and pitiless power! Darwin himself felt all the suffering and
pain that the struggle for existence brings with it. ``We behold the
face of nature bright with gladness, we often see superabundance of
food; we do not see, or we forget, that the birds which are idly
singing round us mostly live on insects or seeds, and are thus
constantly destroying life; or we forget how largely these songsters,
or their eggs, or their nestlings, are destroyed by birds and beasts
of prey; we do not always bear in mind, that though food may be now
superabundant, it is not so at all seasons of each recurring year.''
But he consoles himself with the thought that ``from the war of
nature, from famine and death, the most exalted object which we are
capable of conceiving'' is produced.

It is not easy to see what advantage the common belief in the
``purposiveness of nature'' has over Darwin's doctrine here. For what
a sharp conflict there is between the all-powerful and provident wisdom
that is supposed to have produced everything, and the pain and need,
the suffering to which all living beings are subject---which neither
Darwinists nor anti-Darwinists can overlook! Here is something that
no theory, no worldview can explain. The highest we can attain is to
fix our gaze on the goal achieved through these struggles, and to see
how that pain is a consequence of development toward this goal.
Alongside the assurance that this development goes on unchecked, the
necessity of resignation shows itself; we see the definite,
inexorable conditions without which nothing can be attained here in
the world, and which our wishes and longings cannot in the end change.
But this resignation any worldview whatsoever must necessarily demand,
in one form or another.

Just as the astronomical theory of Kant and Laplace and the geological
theory of Lyell, so Darwin's biological theory teaches us that the
same forces that at every moment operate imperceptibly around us have,
through their constant, never-interrupted operation, produced the
forms in nature that we admire to such a degree that we are inclined
to conceive of them as produced in an instant, by an almighty act of
creation. Just as Copernicus opened up the infinity of the universe
to our gaze, put an end to the old opposition between heaven and
earth, and led us to regard our globe as a vanishing point in the
infinity of worlds, so these more recent theories lead us into the
infinity and eternity of forces. The most decisive oppositions are
only relative stages in the infinite process. We cannot encompass the
infinite; but all the differences our thought seeks to establish
perpetually dissolve anew, prove to be relative. In the infinity of
force and space, the infinity of mind finds its symbol. The highest
Idea we can form is the thought of an infinite world-harmony in which
every individual, in its place and at its time, plays its role and
through its struggle and suffering makes a real contribution to the
great life-process.

When one sees what enthusiasm the Copernican system aroused and how
it gave rise to the first application of the results of natural
science to the outlook on life, one may also expect similar
consequences from these more recent theories---as surely as the
infinity of forces is no less elevating than the infinity of
cosmic space. The inherited doctrines of faith oppose Copernicus and
Newton as they now oppose Lyell and Darwin; but as they had to yield
before the former, so they will also come to make terms with the
latter. We are now simply compelled to form our highest ideas in
accordance with the way in which the universe presents itself to us.
Our outlook on life is a worldview. Just as the discoveries of
Copernicus and Newton compel thought to let its highest Idea encompass
the infinity of the universe---which Giordano Bruno already breathed
in with a conviction and enthusiasm that did not even fail him on the
Inquisition's funeral pyre---so, in accordance with the theory of
development, the conception of the mode of operation of the highest
Idea must now undergo a corresponding change. Not by sudden
intervention, not through unprepared catastrophes and revolutions, but
quietly and imperceptibly operates the eternal force, in whose being
all conditions, all real causes, are elements.

I have heard a talented woman say that if Darwin were right, and we
really descended from the apes, life would lose all worth and
significance for her. This is perhaps the point that stands as most
forbidding for most people. But one who sees that even the highest and
noblest nevertheless has its definite conditions, that it cannot exist
without them and subsists only by appropriating them---one who
recognizes that the apparently absolute differences dissolve in the
infinite connection---will see the matter differently. He will also in
the human mind see a form of the struggling Idea, and will see no
necessity for the spiritual to be breathed into matter from outside or
from above; the spiritual does not stand as an alien element in
existence; it emerges at its determinate stage in development, works
its way forward out of matter itself, and provides in the final
instance the only explanation of matter. The disconcerting quality of
the representation that human beings are supposed to descend from the
apes has its simple ground in the fact that here two developmental
stages are immediately juxtaposed that lie infinitely far apart; what
is essential is thereby skipped over---development itself. How great a
chasm lies between the cannibals from whom we perhaps descend and the
civilized human beings of our age!

We see the mind of the individual human being develop from a modest
beginning; we see in history the human race develop from animal
crudeness to spiritual cultivation and humane social life. What
Darwinism asks is therefore only that one conceive the lines that this
course of development describes as extended in a downward direction,
just as in all our ideal dreams, in our striving for progress and
higher development, we conceive of them as extended in an upward
direction. This leads us to the last point we shall bring forward in
elucidating Darwinism: its ethical character. But before we leave the
domain of the general worldview, I would like, as an
\textit{argumentum ad hominem}, to recall that a Danish philosopher
who stands for all of us as a singularly noble and humane thinker, and
whose outlook on life had a particular warmth and depth of feeling---
that the late F.\ C.\ Sibbern held the same view of the origin of the
human being that one is now so highly offended to encounter in
Darwinist authors. For it says in his work \emph{Speculative
Cosmology}

\noindent (Copenhagen 1846), pp.\ 20--21: ``If the first human being
came into existence in a different way from us, then it is still not
created in any other sense of the word `created' than we are. How
then did the first human being come to be? Without doubt in this way:
that a certain animal formation, from an initial modest beginning,
through a series of generations or offspring---inasmuch as in the
progeny the process of development continually went further and
further---arrived at the point where its last offspring became a being
who, so to speak, opened the eye of spiritual life, saying `I' to
itself and feeling itself in its selfhood, but now also feeling the
drive in itself to name all animals, that is, to impress its stamp
upon them, and to regard itself as standing in relation to them all,
and even above them all. Herewith humanity arose out of animality.''

\section*{III.}

Just as it is commonly held that the metaphysics of Darwinism is
materialism, so it is held that its ethics must amount to animal
sensualism. ``Let us eat and drink, for tomorrow we die!''---this is
the train of thought one generally regards as the natural one for a
Darwinist. Many go still further and hold that suicide is the natural
end for one who entertains such views. ``Every suicide,'' wrote an
American clergyman to Tyndall, ``every suicide in our country (and
they occur daily) is an indirect consequence of the bestial doctrines
that you, Darwin, Spencer, Huxley, et id omne genus, are
preaching.''\footnote{The letter is reprinted in H.\ Spencer's
\textit{The Study of Sociology}, p.\ 419.}

There is scarcely greater proof of the narrow spiritual horizon that
most people still possess than the great difficulty they have in
conceiving that a moral and noble human life can be lived on the basis
of a worldview other than their own. For most people, someone who
believes differently is an immoral being. And behind the accusation
of immorality lies something worse: the hatred of fanaticism. Thus the
American clergyman continues in his letter: ``Shall I not hate those
who hate Thee, Lord!''

In contrast to the passions that the conflict sets in motion, we must
hold to the matter itself and examine it. Darwinism is only a
particular form or application of the universal law of development,
which before Darwin's time had already been recognized as applying
not only in the domain of nature (Kant, Laplace, Lyell) but also in
the domain of history. Already in Bacon of Verulam we find the idea
of the development and progress of humanity. Pascal and Leibniz
expressed this idea more precisely. But it is in the last century that
the first steps toward its realization are taken, by Turgot and
Condorcet, Lessing and Kant. Ethics, which had previously fixed its
gaze only on the development of the individual personality, now had
to come to see that this development can only take place as a member
of the general development of the race. Development of personal life
through participation in cultural-historical work---this is the moral
principle whose source lies in the Idea of development.

Darwinism shares the same ethical motives as the entire worldview
built upon this Idea. The very great role that development plays in
Darwin's hypothesis points in an ethical direction. Rest and stasis
are not ethical; it is work, becoming, struggle, and progress that
matter, and Darwinism leads us directly to this by impressing upon us
that one who does not fight with all his strength will succumb, and
that life itself is a struggle for existence. Viewed immediately, it
is certainly the blind, egoistic drive of self-preservation that comes
to play the greatest role; but it will gradually open up, so to speak,
the more the course of development brings individuals together and
teaches them to conceive of the common higher law to which the
individual wills must bow. The drive of self-preservation is thereby
refined into an ideal striving to preserve and develop the true human
existence both in others and in oneself. Thus the highest ethical
striving remains still a struggle for existence, and the natural law
that Darwin has demonstrated acquires its validity in the ethical
domain as well.

If Darwin were right in his assumption that variation has no limits,
this would open up an even more magnificent and inspiring prospect for
all higher striving. The human being would then not encounter limits
to his ideal longing; what he could not himself attain, he would be
able to lay the foundation of for the generations after him. And
every advance however small, every effort however inconspicuous and
apparently insignificant, would still make its contribution to
development; for it is through an infinite number of infinitely small
operations that the great results are achieved.

On the other hand, Darwinism contains a teaching that the earlier
spiritualist ethics often overlooked: the great significance of
external conditions. The work for the advancement of the human race
consists not least in the improvement of its material circumstances.
Darwinism shows us what significance these circumstances have for the
development of a species---and therefore also for humanity's. It
therefore becomes an ethical task to provide the necessary conditions
of development for all. Darwinism's ethical consequences point thus
in the direction of social reform. It shows us what wildness and what
horror arise when the struggle turns on material self-preservation,
and impresses upon us the necessity of bringing it about that the
struggle---without which human life would stagnate and be extinguished
---can more and more be carried on in a manner worthy of humanity, and
for ends that rise above the blind animal drives, which operate only
because need calls them forth and because they are not forced back by
higher ideas.

Our conclusion, then, is as follows. Darwinism is a hypothesis, and
it is illogical and unscientific to regard it as more. But it does
not merely appeal to a real and significant cause; it has itself grown
out of the necessary presupposition that, in the case of the origin
of species as well, it must have been natural causes that were
operative. The unfortunate theoretical and practical consequences that
have been assumed to follow necessarily from it are unfounded, and
the resistance to Darwinism in fact concerns the entire rational and
humane worldview, which does not stand or fall with Darwinism. This
worldview had been developed in its decisive outlines before Darwinism
appeared, and will continue to develop even if Darwinism should be
refuted. But on the other hand, its truth would provide that worldview
with a significant accretion, enriching it with new perspectives and
ideas. --- But, as has been said: one who knows does not hurry!

\end{document}
